\documentclass[11pt]{article}
\usepackage[margin=0.9in]{geometry}
\usepackage{amsmath,amssymb,amsthm,booktabs}
\usepackage{mathrsfs}
\usepackage[T1]{fontenc}
\usepackage{lmodern,microtype}
\usepackage[colorlinks=true,linkcolor=blue,urlcolor=blue]{hyperref}
\newtheorem{theorem}{Theorem}[section]
\newtheorem{lemma}[theorem]{Lemma}
\newtheorem{proposition}[theorem]{Proposition}
\newtheorem{corollary}[theorem]{Corollary}
\theoremstyle{remark}
\newtheorem{remark}[theorem]{Remark}
\newcommand{\norm}[1]{\left\|#1\right\|}
\title{Toward a four-dimensional Yang--Mills construction\\
\large Notebook 15: a uniform nonlinear constrained Wilson minimum}
\author{Research working document}
\date{16 October 2026}
\begin{document}
\maketitle
\begin{abstract}
We assemble a nonlinear small-field minimizer construction for the
bare Wilson action on homogeneous finite periodic regulators.
The constants and coarse-data analyticity radius are uniform in
the refinement and volume. The proof requires an equal-exponent
constrained Schauder estimate, an exact Wilson force factorization
with removable mesh factors, and a two-Hilbert-direction bound
for the nonlinear constraint curvature. These yield a strict
minimum in a fixed flat gauge slice for the corrected average.
A constraint-preserving gauge flow and a separate transfer to
ordinary group averages then give a minimum, modulo gauge,
among sufficiently small global physical C1-alpha potentials
with prescribed ordinary coarse average. This is a finite-regulator
classical ingredient for a constructive program. It does not yet
construct an interacting continuum quantum theory.
\end{abstract}

\noindent\textbf{Status and provenance.}
Research proofs subjected to separate internal mathematical audits;
no external certification or novelty claim. Source matching refers
to the deposited transcriptions at commit
\texttt{e392ef45013a581664bab76b90867170b30fe10e} of
\href{https://github.com/ethanwsussman/Yang--Mills-LLM-experiment/tree/e392ef45013a581664bab76b90867170b30fe10e}{Yang--Mills-LLM-experiment}.
Original CMP98 and CMP99b pages have not all been independently
recovered. The action and physical pairing match the deposited
four-dimensional bare Wilson conventions; their source-coordinate
Hessian is distinguished explicitly from the ordinary-potential
Hessian used in the nonlinear argument.

\section{The result and the remaining construction problem}
\label{nm15-overview}

Fix a matrix group $SU(N_c)$, normalized trace
$\tau=\operatorname{Tr}/N_c$, an odd block ratio $L\ge3$,
and $0<\alpha<1$. Fine spacing is $\eta=L^{-j}$ and the
retained coarse spacing is one. The regulators are compatible
finite periodic lattices. Potentials are Hermitian traceless
on the real locus and are complexified for analytic estimates.
The complete-field norm $N_\alpha$ controls the potential,
its physical first differences and their $\alpha$-H\"older
seminorms. Transposes always use physical sums. Fixed collar
enlargements and orientation multiplicities are included in
the constants, which may depend on $L,N_c,\alpha$.

The final theorem is Theorem~\ref{ot-minimum}. In words,
for sufficiently small arbitrary coarse data $B$ there is an
analytic small fine field with the exact ordinary homogeneous
constraint $M^{(j)}(U)=e^{iB}$, and its action is minimal among
all fields in a fixed small $N_\alpha$ ball satisfying that
constraint. Equality within the ball means gauge equivalence.
The corrected representative also obeys the sharper expansion
$A_*(B)=HB+O(\norm{B}_\infty^2)$ in $N_\alpha$, where $H$
is the flat minimizing extension. All these bounds are uniform
in the mesh and the compatible periodic volume.

There are three distinct operations in the proof. The flat
Green operator solves the nonlinear stationarity equation in
a fixed gauge slice. A genuine gauge flow covers nearby fields
with the \emph{same corrected constraint}. Finally, two explicit
gauge constructions relate the corrected and ordinary constraint
sets. Neither last step is a consequence of fixed-slice convexity
alone. The ordinary representative need not be in the same flat
gauge slice as the corrected one.

The flat operator identities, canceled Fourier aliases, and
physical exponential locality used below were derived in
Notebooks 8--12. The exact nonlinear constraint chart is from
Notebook 13. Notebook 14 supplies the ordinary/corrected gauge
relation and reference-coordinate comparison. This document
contains the additional estimates that close the nonlinear
minimum argument. Its symbols are restated at their use, but
the earlier flat alias estimates remain explicit prerequisites.

This closes a concrete classical small-field step. It leaves
open extension to local gauges and mixed retained histories,
control of large fields, integration of interacting fluctuations,
running-coupling estimates, and a nontrivial continuum observable
hierarchy. Euclidean invariance, positivity, and the physical
mass gap remain separate obligations. The gap in the constrained
flat auxiliary problem is not a physical Yang--Mills mass gap.


% No preamble. Unique prefix cg-. Requires amsmath/amssymb/amsthm and \norm.
% Independently audited proof module for the fixed-flat-gauge nonlinear minimizer.

\section{The constrained flat inverse preserves the H\"older exponent}
\label{cg-section}

We now establish the constrained Green estimate needed for the
nonlinear fixed-gauge argument. The estimate concerns the
homogeneous flat operators. It is not a nonflat Schauder theorem.
All fine sums use physical measure $\eta^d$; transposes below use
that measure and are bilinear after complexification.

\subsection{Definitions and exact constraints}
\label{cg-algebra}

Set
\begin{align}
 K&=d^Td,&\mathcal L&=D^TD,&
 \mathcal A_0&=K+DRD^T+Q^Ta_{\rm pen}Q,&G_0&=\mathcal A_0^{-1},
                                                        \label{cg-flat}\end{align}
where $D$ is the scalar fine gradient, $d$ the bond curl, and
$Q=Q_M^{\rm st}$ the flat bond average. The one-form penalty is
fixed in a compact positive interval. Define
\begin{equation}\label{cg-extension}
 q_{\rm aux}=(QG_0Q^T)^{-1},\qquad
 H=G_0Q^Tq_{\rm aux},\qquad
 \Gamma=G_0-Hq_{\rm aux}^{-1}H^T.
\end{equation}
The matched flat extension satisfies $QH=I$ and $RD^TH=0$.
With the scalar source projection and its constrained biharmonic
inverse, write
\begin{equation}\label{cg-projection}
 Z=G'RG',\qquad \Lambda=Z\mathcal LD^T=G'RD^T,
 \qquad \Pi=I-D\Lambda.
\end{equation}
The established flat identities give
\begin{equation}\label{cg-scalar-identities}
 Q_0Z=0,\quad \mathcal LZ\mathcal L=R,\quad
 Z\mathcal L^2Z=Z,\quad G_0\Lambda^T=DZ.
\end{equation}
Since $\Lambda H=0$, they imply
$\Gamma\Lambda^T=DZ$,
$\Lambda\Gamma=ZD^T$, and
$\Lambda\Gamma\Lambda^T=Z$.
Expanding the two projection factors therefore proves
\begin{equation}\label{cg-constrained}
 \boxed{\ \mathcal C=\Pi\Gamma\Pi^T=\Gamma-DZD^T\ }.
\end{equation}

The cochain relation $QD=D_cQ_0$ gives $Q\Pi=Q$.
Also $RD^T\Pi=RD^T-R\mathcal LZ\mathcal LD^T=0$.
Consequently
\begin{equation}\label{cg-constraints}
 Q\mathcal C=0,\qquad RD^T\mathcal C=0,
 \qquad \mathcal C Q^T=0.
\end{equation}
The last identity uses symmetry of $\mathcal C$.
Since $\Gamma\mathcal A_0=I-HQ$, any
$\xi\in\mathcal X:=\ker Q\cap\ker(RD^T)$ satisfies
$\Gamma K\xi=\xi$. The identity $D^TK=0$ then gives
\begin{equation}\label{cg-inverse-identities}
 \mathcal CK\xi=\xi\quad(\xi\in\mathcal X),
 \qquad \mathcal CKH=0.
\end{equation}
For the second assertion, use
$KH=Q^T(q_{\rm aux}-a_{\rm pen}I)$ and
\eqref{cg-constraints}. On real regulators
$\mathcal C=\Pi\Gamma\Pi^T$ is nonnegative and is exactly
the inverse of the Maxwell form restricted to $\mathcal X$:
its range lies in $\mathcal X$, and
$\langle\xi,K\mathcal Cf\rangle=\langle\xi,f\rangle$
for every $\xi\in\mathcal X$.
The restricted form is coercive because there it equals the
positive form $\mathcal A_0$.

\subsection{Uniform statement}
\label{cg-statement}

Let $S_{\rm jet}=(I,\nabla)$ contain every ordinary fine first
difference of bond components. Finite component selections,
shifts and curl contractions are allowed. For a fixed
$0<\alpha<1$, use complete-field H\"older norms on fixed unit-cell
collars and admissible positive weights satisfying
$\varpi(n)/\varpi(m)\le e^{\kappa d(n,m)}$.
The norm $N_{\alpha,\varpi}$ contains the field supremum,
first-difference supremum, and the $\alpha$-H\"older seminorm of
the first differences. The input norm
$\norm{f}_{\alpha,\varpi}$ contains the weighted supremum and
the collar $\alpha$-H\"older seminorm.

\begin{theorem}\label{cg-schauder}
For a sufficiently small fixed spatial exponent $\kappa>0$,
uniformly in refinement and compatible periodic volume,
\begin{align}
 N_{\alpha,\varpi}(\mathcal Cf)
       &\le C_\alpha\norm{f}_{\infty,\varpi},
                                                    \label{cg-bounded}\\
 N_{\alpha,\varpi}
   (\mathcal CS_{\rm jet}^T(f_0,f_1))
       &\le C_\alpha\bigl(\norm{f_0}_{\infty,\varpi}
                          +\norm{f_1}_{\alpha,\varpi}\bigr).
                                                    \label{cg-jet}
\end{align}
The same estimates hold on the full lattice. Thus the derivative
datum has the same H\"older exponent as the output gradient.
No assertion is made at either endpoint $\alpha=0,1$.
\end{theorem}

\subsection{The extension subtraction in Gamma}

The established flat equal-exponent estimate gives
$N_{\alpha,\varpi}(G_0\nabla^Tf)
\le C_\alpha\norm{f}_{\alpha,\varpi}$, and the bounded-forcing
estimate gives $G_0:\ell^\infty\to N_\alpha$.
The localized extension satisfies
\begin{equation}\label{cg-H-locality}
 \norm{H(\cdot,b)}_{C^{1,\alpha}(\text{collar }n)}
       \le C_\alpha e^{-\rho d(n,b)}.
\end{equation}
Moreover $q_{\rm aux}^{-1}=QG_0Q^T$ has an exponentially
summable coarse kernel. Hence
\begin{equation}\label{cg-Gamma-derivative}
 \Gamma\nabla^T=G_0\nabla^T
          -Hq_{\rm aux}^{-1}(\nabla H)^T
\end{equation}
maps $C^\alpha_\varpi$ to $N_{\alpha,\varpi}$ uniformly.
In fact its subtracted term accepts bounded data: the physical
transpose row is an integral over cells of volume one, so
\eqref{cg-H-locality} bounds it without any power of the fine-cell
dimension. The weighted exponential kernels can then be convolved.
The undifferentiated argument also proves
$\Gamma:\ell^\infty_\varpi\to N_{\alpha,\varpi}$.

\subsection{A massive biharmonic comparison}
\label{cg-biharmonic}

It remains to estimate $DZD^T\nabla^T$. Put
$Z_m=(\mathcal L+1)^{-2}$. We first prove the stronger remainder
bound
\begin{equation}\label{cg-remainder}
 N_{\alpha,\varpi}
   (D(Z-Z_m)D^T\nabla^Tf)
       \le C_\alpha\norm{f}_{\infty,\varpi}.
\end{equation}
The canceled scalar alias formula is
\begin{align}
 Z_{00}&=R_{s,2}/D_{s,2},&
 Z_{0s}&=-\frac{a_0^-a_s}{P_s^2D_{s,2}},&
 Z_{r0}&=-\frac{a_r^-a_0}{P_r^2D_{s,2}},\nonumber\\
 Z_{rs}&=\frac{\delta_{rs}}{P_r^2}
       -\frac{P_0^2a_r^-a_s}{P_r^2P_s^2D_{s,2}},
 &&r,s\ne0,                                   \label{cg-Z-alias}
\end{align}
where $D_{s,2}=a_0a_0^-+P_0^2\sum_{r\ne0}a_ra_r^-/P_r^2$.
These expressions are uniformly holomorphic on a common enlarged
tube, with all displayed canceled denominators bounded away from
zero. A fixed further shrinking does the same for $P_0+1$;
the high massive denominators are already gapped.

The high diagonal difference is exactly
\begin{equation}\label{cg-high-remainder}
 P_r^{-2}-(P_r+1)^{-2}
 =\int_0^\infty t(1-e^{-t})e^{-tP_r}\,dt,
 \qquad r\ne0.
\end{equation}
There are three external differences in
$D(Z-Z_m)D^T\nabla^T$, and at most one further output
difference and an $\alpha$-increment. The small-time absolute
row bound is therefore
\begin{equation}\label{cg-time-budget}
 t(1-e^{-t})t^{-(4+\alpha)/2}\le t^{-\alpha/2},
 \qquad 0<t\le1,
\end{equation}
which is integrable. The established complex heat-image bounds
apply through these fixed derivative orders by semigroup
splitting. The removed central mode is bounded at small times,
and the high spectral gap controls large times.

For the rank terms, the high output vector $a_r^-/P_r^2$ is
multiplied by at most $C(1+|r|)^{2+\alpha}$.
With $w(r)=\prod_\nu(1+|r_\nu|)^{-1}$, its absolute sum is
bounded by
\[
 C\sum_{j\ge0}(1+j)^{d-1}2^{-j(2-\alpha)}<\infty.
\]
The two input differences leave a row bounded by
$C|a_s|(1+|s|)^{-2}$, whose alias $\ell^2$ norm is uniform.
Central terms are bounded. Physical-cell $L^2$ is dominated by
the supremum without a dimension factor. Reconstruction, periodic
gluing of the complete operator, and contour shifting give the
exponentially local version of \eqref{cg-remainder}.
Only the reconstructed operator is periodized; separated central
chart terms need not individually be periodic.

For the massive part use
\begin{equation}\label{cg-massive-heat}
 Z_m=\int_0^\infty t e^{-t}e^{-t\mathcal L}\,dt.
\end{equation}
The output gradient of $DZ_mD^T\nabla^Tf$ has four fine
differences of the heat kernel. The required next difference
has five. For $k=4,5$ the weighted moment estimate is
\begin{equation}\label{cg-moments}
 \sum_v e^{\kappa|v|_1}|v|_1^\alpha
       |\nabla_\eta^k p_t^\eta(v)|
       \le C_{k,\alpha}t^{(\alpha-k)/2},
       \qquad 0<t\le1.
\end{equation}
For $t\ge\eta^2$, split into $k$ heat factors and use the
weighted first-difference estimate at exponential scale
$\kappa+t^{-1/2}$; the extra exponential pays for the moment.
For $t<\eta^2$, expand the fixed $k$ differences into shifts of
size at most $k\eta$. The dimensionless walk has time below one,
giving $C\eta^{\alpha-k}\le Ct^{(\alpha-k)/2}$.
Thus no sub-mesh time range is omitted.

Every derivative row has sum zero. Subtracting the value of $f$
at the output point and using the complete weighted H\"older
increment bound gives
\begin{align}
 \norm{t\nabla^4e^{-t\mathcal L}f}_{\infty,\varpi}
   &\le Ct^{-1+\alpha/2}\norm{f}_{\alpha,\varpi},\nonumber\\
 \norm{t\nabla^5e^{-t\mathcal L}f}_{\infty,\varpi}
   &\le Ct^{-3/2+\alpha/2}\norm{f}_{\alpha,\varpi}.
                                                        \label{cg-cancellation}
\end{align}
For output separation $\eta\le r\le1$, split at $t=r^2$:
below it use twice the first bound, and above it use $r$ times
the second along a fine path in an enlarged fixed collar.
The gradient difference, divided by the input norm, is at most
\begin{equation}\label{cg-schauder-split}
 C\int_0^{r^2}t^{-1+\alpha/2}\,dt
 +Cr\int_{r^2}^1t^{-3/2+\alpha/2}\,dt+Cr
 \le C\left(\frac2\alpha+\frac2{1-\alpha}+1\right)r^\alpha.
\end{equation}
The last term is the massive large-time tail, with a sufficiently
small fixed exponential weight. Larger separations use the
supremum bound. The gradient supremum is integrable for
$\alpha>0$; the field supremum has absolute small-time budget
$t\,t^{-3/2}=t^{-1/2}$. Therefore
\begin{equation}\label{cg-massive-result}
 N_{\alpha,\varpi}(DZ_mD^T\nabla^Tf)
       \le C_\alpha\norm{f}_{\alpha,\varpi}.
\end{equation}
Combining \eqref{cg-remainder} and \eqref{cg-massive-result}
proves the differentiated subtraction estimate. Without the
last input derivative, the absolute budget is
$t\,t^{-(3+\alpha)/2}$, integrable for $\alpha<1$; this gives
$DZD^T:\ell^\infty_\varpi\to N_{\alpha,\varpi}$.
Equation \eqref{cg-constrained} and
\eqref{cg-Gamma-derivative} now prove
Theorem~\ref{cg-schauder}.

\subsection{Full-lattice identities and the application boundary}

The massive estimates use complete fields and global heat-row
cancellation. The remainder estimates use exponentially local
pointwise cell kernels with input derivatives already on the
kernel. No sharp restriction of a H\"older field to one cell is
used. Exact periodization preserves the estimates, including
small compatible tori and their incidence multiplicities.

At each fixed spacing the identities above extend from Hilbert
fields to bounded or admissibly weighted complete fields by
absolute summability of the operator products. In particular
$R=\mathcal LZ\mathcal L$ is exponentially summable; fixed-mesh
equivalence of cell norms suffices for this identity argument.
It is not necessary to call the gauge projection local.
For $\xi\in\mathcal X\cap N_\alpha$, $K\xi$ can be read as
a first-divergence datum with H\"older coefficients, and
\eqref{cg-inverse-identities} agrees with that interpretation.

This supplies the bounded-forcing and equal-exponent jet-forcing
Green inputs for the fixed-flat-gauge nonlinear stationarity
argument. The nonlinear action force, nonlinear constraint
derivatives and strict-minimum argument are separate ingredients.
No moving-gauge nonflat Schauder estimate, mixed retained-history
result or interacting continuum limit is asserted here.

% Notebook 15 force fragment; no preamble. Unique label prefix wf-.
% Requires standard theorem environments and a one-argument \norm macro.

\section{A uniform nonlinear Wilson force in physical coordinates}
\label{wf-section}

We work first on finite compatible periodic regulators, with
fine spacing $0<\eta\le1$, physical measure $\eta^d$ and
normalized trace $\tau=\operatorname{Tr}/N_c$. The Lie algebra
consists of Hermitian traceless matrices and is complexified for
analytic arguments. All transposes use the fixed bilinear trace
pairing. Constants are independent of refinement and volume.
On the infinite lattice the formulas concern local variations;
an arbitrary bounded field need not have finite total action.

In dimension four the action, trace and physical pairing below
match the deposited conventions: Raw I:454--467 sums the
normalized Wilson density once per unoriented plaquette,
I:474--479 places the coupling outside that action, and
I:12756--12760 specifies the physical bilinear pairing.
The source fluctuation chart is explicit at I:15643--15648.
Here we differentiate the same action in the ordinary global
coordinate $U=e^{i\eta A}$; its Hessian at nonzero $A$ is not
silently identified with a differently based left-exponential
Hessian. The original journal texts have not all been
independently reverified.

\subsection{Exact plaquette density and its removable scale factor}

For $p=(x;\mu,\nu)$, $\mu<\nu$, define the physical jet
\begin{equation}\label{wf-jet}
 (\mathsf S_\eta A)_p=(a,b,c,e)
  =\bigl(A_\mu(x),A_\nu(x),
       \nabla_{\eta,\mu}A_\nu(x),
       \nabla_{\eta,\nu}A_\mu(x)\bigr).
\end{equation}
Here $\nabla_{\eta,\rho}f(x)
=\eta^{-1}(f(x+\eta e_\rho)-f(x))$. The plaquette
holonomy is exactly
\begin{equation}\label{wf-plaquette}
 P_\eta(a,b,c,e)
  =e^{i\eta a}e^{i\eta b+i\eta^2c}
                    e^{-i\eta a-i\eta^2e}e^{-i\eta b}.
\end{equation}
This merely rewrites shifted link potentials; it is not a
Taylor approximation. In potential coordinates set
\begin{equation}\label{wf-action}
 \mathcal W_\eta(A):=\mathcal S_\eta(e^{i\eta A})
   =\eta^{d-4}\sum_p
          \left[1-\frac12\tau(P_p+P_p^{-1})\right].
\end{equation}
In $d=4$ the prefactor is one, as in the deposit. For other
fixed $d$ the displayed normalization defines the model used
in this local calculation.

\begin{lemma}\label{wf-removable}
The matrix
\[
 \mathfrak F_\eta(z)=\frac{P_\eta(z)-I}{\eta^2}
\]
extends jointly holomorphically through $\eta=0$. The density
\begin{equation}\label{wf-density}
 L_\eta(z)=-\frac12\tau\bigl(\mathfrak F_\eta(z)^2
                                              P_\eta(z)^{-1}\bigr)
\end{equation}
has uniformly bounded jet derivatives of every fixed order
for $\eta\in[0,1]$ on a fixed smaller complex jet ball. Its
quadratic Taylor part at $z=0$ is exactly
\begin{equation}\label{wf-quadratic}
 L_2(z)=\frac12\tau((c-e)^2),
\end{equation}
independently of $\eta$. Moreover
\begin{equation}\label{wf-exact-action}
 \mathcal W_\eta(A)=\eta^d\sum_p
                           L_\eta((\mathsf S_\eta A)_p).
\end{equation}
\end{lemma}

\begin{proof}
The matrix product \eqref{wf-plaquette} is entire in $\eta$
and the finitely many complex jet coordinates. For every $z$,
$P_0(z)=I$ and
$\partial_\eta P_\eta(z)|_{\eta=0}=i(a+b-a-b)=0$.
Thus the extension is explicitly
\begin{equation}\label{wf-removable-integral}
 \mathfrak F_\eta(z)
  =\int_0^1(1-s)\,\partial_\eta^2P_{s\eta}(z)\,ds,
\end{equation}
where the derivative is with respect to the scalar argument
before evaluation at $s\eta$. The inverse $P_\eta^{-1}$ is
the reverse product of inverse exponentials and is entire.
On $|\eta|\le2$ and a fixed bounded jet polydisc these
expressions are bounded. Cauchy estimates in the jet
coordinates give the asserted uniform derivative bounds
on a smaller polydisc.

At $z=0$, $\mathfrak F_\eta=0$, and direct differentiation
gives
\begin{equation}\label{wf-linear-F}
 D_zP_\eta(0)[z]=i\eta^2(c-e),\qquad
 D_z\mathfrak F_\eta(0)[z]=i(c-e).
\end{equation}
Equation \eqref{wf-quadratic} follows from
\eqref{wf-density}. In particular
$\mathfrak F_0(z)=i(c-e)-[a,b]$; the commutator is quadratic
in the jet and does not change the quadratic density.
Finally, the exact algebraic identity
\[
 P+P^{-1}-2I=(P-I)^2P^{-1}
\]
proves \eqref{wf-exact-action}. All inverse mesh factors
cancel before a force derivative is taken.
\end{proof}

\subsection{The exact force and its local nonlinear coefficient}

Put $R_\eta=L_\eta-L_2$ and define the jet gradient $f_\eta$
by the fixed bilinear trace pairing:
\begin{equation}\label{wf-gradient-definition}
 DR_\eta(z)[w]=\sum_{\ell\in\{a,b,c,e\}}
                                \tau(f_{\eta,\ell}(z)w_\ell).
\end{equation}
Projection onto the fixed traceless space is included. On a
fixed small complex jet ball, uniformly in $\eta$,
\begin{equation}\label{wf-pointwise-bounds}
 \begin{gathered}
 f_\eta(0)=Df_\eta(0)=0,\qquad
 \norm{f_\eta(z)}\le C\norm{z}^2,\qquad
 \norm{Df_\eta(z)}\le C\norm{z},\\
 \norm{D^2f_\eta(z)}+\norm{D^3f_\eta(z)}\le C.
 \end{gathered}
\end{equation}
These follow from Lemma~\ref{wf-removable} and Taylor's
formula, since $R_\eta$ begins at degree three.

Let $d$ be the ordinary fine curl, so $(dA)_p=c-e$, and
let $K=d^Td$. Define $E_\eta(A)$ by
$D\mathcal W_\eta(A)[H]=\langle E_\eta(A),H\rangle_\eta$.
Then
\begin{equation}\label{wf-force-factorization}
 E_\eta(A)=KA+\mathsf S_\eta^T\mathcal N_\eta(A),
 \qquad
 \mathcal N_\eta(A)_p=f_\eta((\mathsf S_\eta A)_p).
\end{equation}
For a four-slot plaquette field $q$, the transpose is
explicitly
\begin{equation}\label{wf-jet-transpose}
 \begin{split}
 (\mathsf S_\eta^Tq)_\lambda(x)
 ={}&\sum_{\nu>\lambda}\left[
       q_{\lambda\nu}^{a}(x)
                +\nabla_{\eta,\nu}^Tq_{\lambda\nu}^{e}(x)\right]\\
 &+\sum_{\mu<\lambda}\left[
       q_{\mu\lambda}^{b}(x)
                +\nabla_{\eta,\mu}^Tq_{\mu\lambda}^{c}(x)\right],
 \end{split}
\end{equation}
where $\nabla_{\eta,\rho}^Tq(x)
=\eta^{-1}(q(x-\eta e_\rho)-q(x))$. Every slot has the same
physical measure $\eta^d$. At quadratic order the $c,e$
coefficients are $c-e,e-c$, confirming $K=d^Td$ with
coefficient one. The derivative in
\eqref{wf-jet-transpose} stays outside the nonlinear
coefficient; a uniform pointwise bound on the unprocessed
nonlinear force is not asserted.

\subsection{Physical H\"older norms without a mesh loss}

Fix $0<\alpha<1$. All fields are complete fields, with
componentwise comparisons at translated sites. Define
\begin{equation}\label{wf-Nalpha}
 \begin{split}
 N(A)&=\norm{A}_\infty+
                  \max_\rho\norm{\nabla_{\eta,\rho}A}_\infty,\\
 N_\alpha(A)&=N(A)+
                  \max_\rho[\nabla_{\eta,\rho}A]_\alpha.
 \end{split}
\end{equation}
The seminorms use physical distances in fixed enlarged
unit-cell collars. Suprema supply longer-distance
increments. This is equivalent to the complete-field
$C^{1,\alpha}$ norm used in the minimizer argument, with
constants independent of $\eta$.

For two points at physical distance at most a fixed $R$,
telescoping along an axis-aligned path gives
\begin{equation}\label{wf-holder-telescoping}
 \frac{|A(x)-A(y)|}{|x-y|^\alpha}
 \le\sqrt d\,R^{1-\alpha}
                      \max_\rho\norm{\nabla_{\eta,\rho}A}_\infty.
\end{equation}
At nearest-neighbor separation the factor is
$\eta^{1-\alpha}$, not $\eta^{-\alpha}$. The derivative
jet slots already have their H\"older seminorms in
\eqref{wf-Nalpha}. Consequently
\begin{equation}\label{wf-jet-norm}
 \norm{\mathsf S_\eta A}_\infty\le C_JN(A),\qquad
 \norm{\mathsf S_\eta A}_{C^\alpha}\le C_JN_\alpha(A).
\end{equation}
No second-difference supremum is needed.

For a positive cell weight $\varpi$ with
$\varpi(n)/\varpi(m)\le e^{\kappa d(n,m)}$, weight each
complete local collar norm by $\varpi(n)$. Denote the
resulting norms by $C^\alpha_\varpi$, $N_\varpi$ and
$N_{\alpha,\varpi}$. Neighboring fixed collars have
comparable weights. The same telescoping argument gives
\begin{equation}\label{wf-weighted-jet}
 \norm{\mathsf S_\eta H}_{C^\alpha_\varpi}
       \le C_{J,\kappa}N_{\alpha,\varpi}(H),\qquad
 \norm{\mathsf S_\eta H}_{\infty,\varpi}
       \le C_{J,\kappa}N_\varpi(H).
\end{equation}
No derivative is taken of a weight and no field is cut
sharply into cells. Periodic lifting gives the same bounds
on compatible tori.

\begin{theorem}\label{wf-nonlinear-theorem}
On a fixed small complex $N_\alpha$ ball, $\mathcal N_\eta$
is holomorphic and, uniformly in refinement and volume,
\begin{align}
 \norm{\mathcal N_\eta(A)}_{C^\alpha_\varpi}
       &\le C N(A)N_{\alpha,\varpi}(A),
                                           \label{wf-quadratic-bound}\\
 \norm{\mathcal N_\eta(A)-\mathcal N_\eta(B)}_{C^\alpha_\varpi}
       &\le C\bigl(N_\alpha(A)+N_\alpha(B)\bigr)
                             N_{\alpha,\varpi}(A-B),
                                           \label{wf-Lipschitz}\\
 \norm{D\mathcal N_\eta(A)[H]}_{C^\alpha_\varpi}
       &\le C\left[N(A)N_{\alpha,\varpi}(H)
                         +N_\alpha(A)N_\varpi(H)\right].
                                           \label{wf-directional}
\end{align}
Weighted estimates require only that their right sides be
finite. The base fields are in the unweighted small ball;
one may put the weight on a direction or difference.
\end{theorem}

\begin{proof}
On each collar the product identity gives
\[
 [uv]_\alpha\le\norm{u}_\infty[v]_\alpha
                         +[u]_\alpha\norm{v}_\infty.
\]
It remains valid with the weight on either factor.
Write $z=\mathsf S_\eta A$. The supremum bound in
\eqref{wf-pointwise-bounds} and integration of $Df_\eta$
along the segment from $z(x)$ to $z(y)$ give
\[
 |f_\eta(z(x))-f_\eta(z(y))|
       \le C\norm{z}_\infty|z(x)-z(y)|.
\]
Together with \eqref{wf-jet-norm} and
\eqref{wf-weighted-jet}, this proves
\eqref{wf-quadratic-bound}.

The exact derivative is
$Df_\eta(z)\mathsf S_\eta H$. Its coefficient has supremum
at most $C\norm{z}_\infty$ and H\"older seminorm at most
$C[z]_\alpha$, since $D^2f_\eta$ is bounded. The product
estimate proves \eqref{wf-directional}. Integrating it
along $B+s(A-B)$ proves \eqref{wf-Lipschitz}. In particular
the unweighted local Lipschitz constant on a radius-$r$
ball is at most $Cr$.

The uniform local Taylor series for $f_\eta$ starts in
degree two. The H\"older product inequality and
\eqref{wf-jet-norm} make its substitution converge
absolutely on a fixed smaller $N_\alpha$ ball. This proves
holomorphy in the complete Banach norms. Higher derivatives
follow from the same series or Cauchy's formula. With one
weighted direction, use the weighted product inequality
term by term. If a weight is not bounded below, the weighted
statement is read on the intersection with the unweighted
small ball; no embedding into that ball is presumed.
\end{proof}

For example the same proof gives
\begin{equation}\label{wf-second-directional}
 \norm{D^2\mathcal N_\eta(A)[H,K]}_{C^\alpha_\varpi}
 \le C(1+N_\alpha(A))N_\alpha(H)N_{\alpha,\varpi}(K),
\end{equation}
with the weighted direction exchangeable by symmetry.
Uniform bounds on a complex ball give factorial Cauchy
bounds on every smaller ball. All these statements are
about nonlinear coefficients before the exterior
transpose in \eqref{wf-force-factorization}.

\subsection{Energy Hessian and the minimizer interface}

Differentiating \eqref{wf-force-factorization} once gives
\begin{equation}\label{wf-energy-Hessian}
 DE_\eta(A)=K+\mathsf S_\eta^T
                Df_\eta(\mathsf S_\eta A)\mathsf S_\eta.
\end{equation}
The middle coefficient is symmetric for the fixed bilinear
transpose, is local in the plaquette jet, and has supremum
at most $CN(A)$. Thus for physical Hilbert directions $H,K$,
\begin{equation}\label{wf-energy-form}
 \left|D^2\mathcal W_\eta(A)[H,K]
                         -\langle dH,dK\rangle_\eta\right|
 \le CN(A)\norm{\mathsf S_\eta H}_{2,\eta}
                         \norm{\mathsf S_\eta K}_{2,\eta}.
\end{equation}
This energy statement needs small $N(A)$, not a second
background derivative. It is an ordinary-potential
Hessian identity for the exactly specified Wilson action.

For the fixed-gauge minimizer, the equal-exponent estimate
for its selected flat constrained inverse $\mathcal C$,
\[
 \mathcal C\mathsf S_\eta^T:C^\alpha\longrightarrow N_\alpha,
\]
combines directly with Theorem~\ref{wf-nonlinear-theorem}.
The resulting nonlinear map is quadratic and has
Lipschitz constant $O(r)$ on a radius-$r$ ball, also for
a single weighted direction. The constrained Green estimate
is proved separately; a positive loss between input and
output exponents would not suffice for this step. No
additional fine derivative of $\mathcal N_\eta(A)$ or
inverse power of $\eta$ is used in the contraction.

% No preamble. Notebook 15 minimizer module. Unique labels: nm-.
% Requires amsmath, amssymb, amsthm and a one-argument \norm macro.

\section{A uniform local constrained Wilson minimum on the flat gauge slice}
\label{nm-section}

Fix an odd $L\ge3$, $M=L^j$, $\eta=M^{-1}$, final coarse scale one,
and $0<\alpha<1$. The primary statement is on finite compatible
periodic regulators. Constants may depend on $L$, dimension, the
color norm, and $\alpha$, but not on $j$ or periodic volume. All
transposes use the physical fine pairing $\eta^d\sum\tau(AB)$ and
the corresponding coarse pairing; at complex fields they are
bilinear transposes. In this module
$\mathcal S(A)$ means the ordinary-potential action
$\mathcal W_\eta(A)=\mathcal S_\eta(e^{i\eta A})$ of the force
section, and $S_{\rm jet}=\mathsf S_\eta$. Write the average in its
fixed-reference potential coordinate as
\begin{equation}\label{nm-average}
 \mathcal F(A)=QA+\mathcal R(A),\qquad Q=Q_M^{\rm st},\qquad
 \mathcal R(0)=D\mathcal R(0)=0.
\end{equation}
The symbol $R$ without calligraphic font continues to denote the
fixed flat scalar gauge projection. In particular, $R$ and
$\mathcal R$ are different objects.

\subsection{Two uniform estimates for the nonlinear constraint}
\label{nm-average-estimates}

The previously established nonlinear averaging ball gives a
one-step dimensionless corrected map $F_L$ and physical map
\begin{equation}\label{nm-one-step}
 \mathcal F_h(A)=(Lh)^{-1}F_L(hA)=Q_hA+\mathcal R_h(A).
\end{equation}
Let $m_L$ bound the number of input bonds read by one output, and
$N_L$ the number of outputs reading one input. On a fixed local
complex neighborhood choose $K_0$ so that each matrix entry of
$D(F_L-LQ_h)(X)$ has norm at most $K_0\norm{X}_\infty$.
The corresponding remainder entries in physical coordinates are
bounded by $K_0h\norm{A}_\infty/L$. With
\begin{equation}\label{nm-local-constants}
 D_\infty=\frac{m_LK_0}{L},\qquad
 D_2=K_0\sqrt{m_LN_L}\,L^{d/2-1},\qquad
 D_*=N_LK_0L^{d-1},\qquad S=\frac1{L-1},
\end{equation}
the physical transpose therefore satisfies
\begin{equation}\label{nm-adjoint-one-step}
 \norm{D\mathcal R_h(A)^T}_{\infty\to\infty}
 \le D_*h\norm{A}_\infty.
\end{equation}
Indeed the transpose converts an ordinary column sum to a physical
row sum with factor $(Lh)^d/h^d=L^d$. This factor is already
included in $D_*$.

Both $Q_h$ and its physical transpose are bounded-field
contractions. If $a=\norm{A}_\infty$ belongs to the common nonlinear
ball, the intermediate fields obey $\norm{A_t}_\infty\le2a$.
The spacings $h_t=L^{t-j}$ satisfy $\sum_{t<j}h_t\le S$.
Expanding the difference of the transposed nonlinear derivative
product and its flat product thus proves
\begin{equation}\label{nm-adjoint-composite}
 \norm{D\mathcal R(A)^T}_{\infty\to\infty}
 \le e^{2D_*aS}-1\le c_*a
\end{equation}
on a smaller fixed ball. The analogous derivative bounds in the
bounded and physical Hilbert norms follow with $D_\infty$ and
$D_2$. These are absolute physical row or column estimates, not
inferences from the fine dimension. The local analytic expansions
also make $A\mapsto D\mathcal R(A)^T$ holomorphic in this operator
norm. Cauchy's estimate on a slightly larger common ball gives
\begin{equation}\label{nm-adjoint-lipschitz}
 \norm{D\mathcal R(A)^T-D\mathcal R(A')^T}_{\infty\to\infty}
 \le c_{*,1}\norm{A-A'}_\infty.
\end{equation}
No bounded-field estimate for the gauge projection $\Pi$ is used.

We also need a bound with two Hilbert directions, rather than just
two bounded directions. Let $K_2$ bound the local bilinear Hessian
$D^2F_L$ in the maximum norms of its inputs; it bounds the full
second derivative, including its factor two. The scaling
$D^2\mathcal F_h=(h/L)D^2F_L(hA)$, local Cauchy--Schwarz, and
incidence counting give
\begin{equation}\label{nm-hessian-one-step}
 \left|\left\langle w,D^2\mathcal F_h(A)[B,C]\right\rangle_{Lh}\right|
 \le k_2h\norm{w}_\infty\norm{B}_{2,h}\norm{C}_{2,h},
 \qquad k_2=K_2N_LL^{d-1}.
\end{equation}
For clarity, replace each local input maximum by its local
Euclidean norm. Summing the product over output bonds and applying
Cauchy--Schwarz costs at most $N_L$; restoring the fine measure
then converts $(Lh)^d(h/L)$ to the factor in
\eqref{nm-hessian-one-step}.

Let $P_t$ be the derivative of the nonlinear prefix through stage
$t$, and $T_t$ the derivative of the suffix after stage $t+1$,
both evaluated along the actual trajectory. The exact second-order
chain rule is
\begin{equation}\label{nm-hessian-chain}
 D^2\mathcal F(A)[B,C]
 =\sum_{t<j}T_tD^2\mathcal F_{h_t}(A_t)[P_tB,P_tC].
\end{equation}
The two prefix Hilbert norms are each at most $e^{2D_2aS}$ and the
suffix transpose bounded-field norm at most $e^{2D_*aS}$.
Pairing \eqref{nm-hessian-chain} with $w$ and using
\eqref{nm-hessian-one-step} therefore proves the uniform form bound
\begin{equation}\label{nm-hessian-composite}
 \boxed{\quad
 \left|\left\langle w,D^2\mathcal F(A)[B,C]\right\rangle_1\right|
 \le k_2S e^{(4D_2+2D_*)aS}
       \norm{w}_\infty\norm{B}_{2,\eta}\norm{C}_{2,\eta}.
 \quad}
\end{equation}

\subsection{The exact reduced equation}
\label{nm-reduction}

Let $K=d^Td$, let $q=q_M$ be the Maxwell coarse precision, and let
$H$ be the flat minimizing extension in the matched source gauge.
The flat identities are
\begin{equation}\label{nm-flat-identities}
 QH=I,\qquad KH=Q^Tq,\qquad RD^TH=0.
\end{equation}
Set
\begin{equation}\label{nm-slice-green}
 \mathcal X=\ker Q\cap\ker(RD^T),\qquad
 \mathcal C=\Pi\Gamma\Pi^T=\Gamma-DZD^T,
 \qquad \Gamma=G_0-H(q+a_{\rm pen})^{-1}H^T.
\end{equation}
The established constrained-inverse identities give
\begin{equation}\label{nm-constrained-identities}
 Q\mathcal C=RD^T\mathcal C=\mathcal C Q^T=0,\qquad
 \mathcal C K\xi=\xi\quad(\xi\in\mathcal X),\qquad
 \mathcal C KH=0.
\end{equation}
On a real finite regulator, $\mathcal C$ is the positive
variational inverse of $K$ on $\mathcal X$. Also
$H^TK\xi=qQ\xi=0$. The semidefinite $q$ in these identities is
not the positive auxiliary precision $q+a_{\rm pen}$.

For coarse $B$ the exact local constraint chart is
\begin{equation}\label{nm-chart}
 A=\xi+H[B-\mathcal R(A)],\qquad \xi\in\mathcal X.
\end{equation}
It gives $\mathcal F(A)=B$ and $RD^TA=0$. Conversely, every small
field satisfying these two conditions belongs to this chart,
by taking $\xi=A-H[B-\mathcal R(A)]$. Its derivative in $\xi$ is
$T_A=(I+HD\mathcal R(A))^{-1}$ restricted to $\mathcal X$.

Write the ordinary-potential Wilson Euler derivative as
\begin{equation}\label{nm-force}
 E(A)=KA+N(A),\qquad N(A)=S_{\rm jet}^Tf_\eta(S_{\rm jet}A),
\end{equation}
where $S_{\rm jet}$ collects the finitely many local values and
ordinary first differences used by the plaquette jet. Stationarity
of the action in the chart is equivalent to
$\mathcal C T_A^TE(A)=0$. The ordered inverse identity
\[
 T_A^T=I-D\mathcal R(A)^T
       (I+H^TD\mathcal R(A)^T)^{-1}H^T
\]
and \eqref{nm-flat-identities}--\eqref{nm-chart} transform this
condition exactly into
\begin{equation}\label{nm-fixed-point}
 \boxed{\quad
 \xi=-\mathcal C N(A)
 +\mathcal C D\mathcal R(A)^T
       (I+H^TD\mathcal R(A)^T)^{-1}
       \{q[B-\mathcal R(A)]+H^TN(A)\}.
 \quad}
\end{equation}
In particular, the nonlinear constraint contributes the second
term in \eqref{nm-fixed-point}; it cannot be omitted.

\subsection{Uniform analytic contraction}
\label{nm-contraction}

Use the complete-field physical norm
\begin{equation}\label{nm-norm}
 \norm{A}_{\mathcal N_\alpha}
 =\norm{A}_\infty+\norm{\nabla A}_\infty+[\nabla A]_\alpha.
\end{equation}
It controls the complete $C^\alpha$ norms of the jet, uniformly in
$\eta$. Telescoping controls differences of $A$ at distances below
one; its bounded norm handles larger distances. The flat extension
and the constrained Green theorem proved above supply
\begin{align}
 H &: \ell^\infty_c\longrightarrow\mathcal N_\alpha,&
 \mathcal C &: \ell^\infty_f\longrightarrow\mathcal N_\alpha,
                                                        \label{nm-analytic-input-a}\\
 \mathcal C S_{\rm jet}^T &: C^\alpha_f\longrightarrow\mathcal N_\alpha,
 &q &: \ell^\infty_c\longrightarrow\ell^\infty_c.
                                                        \label{nm-analytic-input-b}
\end{align}
The equal exponent in \eqref{nm-analytic-input-b} is essential.
It uses the separate constrained biharmonic cancellation, including
the fourth and fifth heat moments and the separation-dependent
time split, not merely an estimate for $G_0$ or a theorem with
loss of H\"older exponent.

The transpose estimates for $H$ and its first differences give
\begin{equation}\label{nm-transpose-input}
 H^T:\ell^\infty_f\longrightarrow\ell^\infty_c,
 \qquad H^TS_{\rm jet}^T=(S_{\rm jet}H)^T:
       C^\alpha_f\longrightarrow\ell^\infty_c.
\end{equation}
For the second map the bounded input norm already suffices:
the physical volume of a fine cell is fixed and the physical-cell
Hilbert transpose rows of $S_{\rm jet}H$ are summable. Thus
\eqref{nm-transpose-input} never differentiates a rough force
coefficient.

The Wilson jet calculation supplies a holomorphic function
$f_\eta$ with $f_\eta(0)=Df_\eta(0)=0$ and the uniform estimates
\begin{align}
 \norm{f_\eta(S_{\rm jet}A)}_{C^\alpha}
 &\le c_f\norm{A}_{\mathcal N_\alpha}^2,                 \label{nm-force-bound}\\
 \norm{f_\eta(S_{\rm jet}A)-f_\eta(S_{\rm jet}A')}_{C^\alpha}
 &\le c_f(\norm{A}_{\mathcal N_\alpha}+
           \norm{A'}_{\mathcal N_\alpha})
                 \norm{A-A'}_{\mathcal N_\alpha}.        \label{nm-force-lipschitz}
\end{align}
These estimates concern the coefficient before the outside
transpose in \eqref{nm-force}.

Shrink the common ball so that
\begin{equation}\label{nm-neumann-smallness}
 \norm{H}_{\infty\to\mathcal N_\alpha}
       \norm{D\mathcal R(A)}_{\infty\to\infty}<\frac12,
 \qquad
 \norm{H^T}_{\infty\to\infty}
       \norm{D\mathcal R(A)^T}_{\infty\to\infty}<\frac12.
\end{equation}
The first inequality makes \eqref{nm-chart} an analytic
$\mathcal N_\alpha$ chart. Put $b=\norm{B}_\infty$ and
$h_N=\norm{H}_{\infty\to\mathcal N_\alpha}$. Its contraction
estimate gives
\begin{equation}\label{nm-chart-estimates}
 \norm{A(\xi,B)}_{\mathcal N_\alpha}
 \le2(\norm{\xi}_{\mathcal N_\alpha}+h_Nb),\qquad
 \norm{A(\xi,B)-A(\xi',B)}_{\mathcal N_\alpha}
 \le2\norm{\xi-\xi'}_{\mathcal N_\alpha}.
\end{equation}
Define the coarse multiplier
\begin{equation}\label{nm-multiplier}
 m(A,B)=(I+H^TD\mathcal R(A)^T)^{-1}
             \{q[B-\mathcal R(A)]+H^TN(A)\}.
\end{equation}
For $\norm{A}_{\mathcal N_\alpha},\norm{A'}_{\mathcal N_\alpha}\le r$
in this chart, the preceding bounds imply
\begin{equation}\label{nm-multiplier-bounds}
 \norm{m(A,B)}_\infty\le C(b+r^2),\qquad
 \norm{m(A,B)-m(A',B)}_\infty
 \le C(b+r)\norm{A-A'}_{\mathcal N_\alpha}.
\end{equation}
For the first bound use inverse norm at most two,
$\norm{\mathcal R(A)}_\infty\le Cr^2$, and
$H^TN=(S_{\rm jet}H)^Tf_\eta$. For the second use the inverse
difference identity and \eqref{nm-adjoint-lipschitz}; the change
of the braces is $O(r\norm{A-A'})$, and the change of the inverse
multiplies a quantity $O(b+r^2)$.

Let $\Phi_B(\xi)$ denote the right side of
\eqref{nm-fixed-point}, with $A=A(\xi,B)$. Then
\begin{align}
 \norm{\Phi_B(\xi)}_{\mathcal N_\alpha}
 &\le C(r^2+rb+r^3),                                   \label{nm-contraction-size}\\
 \norm{\Phi_B(\xi)-\Phi_B(\xi')}_{\mathcal N_\alpha}
 &\le C(b+r)\norm{\xi-\xi'}_{\mathcal N_\alpha}.          \label{nm-contraction-lipschitz}
\end{align}
To check the constraint contribution, expand the difference of
$D\mathcal R(A)^Tm(A,B)$ into its two factors and apply
\eqref{nm-adjoint-composite}, \eqref{nm-adjoint-lipschitz}, and
\eqref{nm-multiplier-bounds}; then use the bounded-input map
$\mathcal C$ in \eqref{nm-analytic-input-a}. The force contribution
uses \eqref{nm-analytic-input-b} and
\eqref{nm-force-lipschitz}. Both contributions lie in $\mathcal X$.

Choose a fixed small $\mathcal X$-ball and then $b_*>0$ so that its
chart image stays in the preceding balls, \eqref{nm-contraction-size}
maps it into itself, and \eqref{nm-contraction-lipschitz} is below
one half. This constructs a unique small solution $\xi(B)$ for
every complex $\norm{B}_\infty<b_*$. Uniform analytic contraction
makes $\xi(B)$ and $A(B)$ holomorphic. The sharper bounds are
\begin{equation}\label{nm-quadratic-remainder}
 \norm{\xi(B)}_{\mathcal N_\alpha}\le Cb^2,
 \qquad \norm{A(B)-HB}_{\mathcal N_\alpha}\le Cb^2,
 \qquad \norm{A(B)}_{\mathcal N_\alpha}\le Cb.
\end{equation}
Indeed on $\norm{\xi}\le Kb^2$, after reducing $b_*$,
\eqref{nm-chart-estimates} gives $r\le4h_Nb$. The right side of
\eqref{nm-contraction-size} is then $C'b^2$. Take $K>2C'$, shrink
$b_*$ again so that $Kb^2\le h_Nb$ and the Lipschitz bound is below
one half, and use uniqueness in the original ball. The chart
remainder $-H\mathcal R(A)$ proves the second estimate.
If a complex coarse ball has distance $\delta$ from the boundary
of the holomorphy ball and $\norm{A(B)}_{\mathcal N_\alpha}\le C_0$
there, multivariable Cauchy gives, for $n\ge1$,
\begin{equation}\label{nm-cauchy}
 \norm{D^nA(B)[B_1,\ldots,B_n]}_{\mathcal N_\alpha}
 \le C_0 n!\left(\frac e\delta\right)^n
                    \prod_{i=1}^n\norm{B_i}_\infty.
\end{equation}
Use polyradii approaching $\delta/n$ and $n^n\le e^n n!$.

\subsection{The real solution is a strict local minimum}
\label{nm-minimum}

For real fields the ordinary-coordinate Wilson Hessian has the
form
\begin{equation}\label{nm-action-hessian}
 D^2\mathcal S(A)=K+S_{\rm jet}^TW(A)S_{\rm jet},
 \qquad \norm{W(A)}_\infty\le C\norm{A}_{\mathcal N_\alpha}.
\end{equation}
This follows from the exact jet density whose quadratic Taylor
part is the Maxwell density. It does not identify two different
coordinate Hessians at a nonzero background.

The flat full inverse's first-derivative energy bound gives
\begin{equation}\label{nm-coercivity}
 \norm{z}_{2,\eta}^2+\norm{\nabla z}_{2,\eta}^2
 \le C_E\langle z,Kz\rangle_\eta\qquad(z\in\mathcal X).
\end{equation}
In fact both the gauge and averaging penalties in the positive
full precision vanish on $\mathcal X$. Apply the uniform bound on
$(I,\nabla)G_0(I,\nabla)^T$ to its square-root factorization.

Let $z\in\mathcal X$ and $v=T_Az$. Differentiating the chart gives
$v=z-HD\mathcal R(A)v$. The cross terms in its Maxwell energy
vanish because $H^TKz=0$, while $q$ is positive semidefinite.
Consequently
\begin{equation}\label{nm-tangent-energy}
 \langle v,Kv\rangle
 =\langle z,Kz\rangle+
   \langle D\mathcal R(A)v,qD\mathcal R(A)v\rangle
 \ge\langle z,Kz\rangle.
\end{equation}
The uniform physical Hilbert bound $\norm{D\mathcal R(A)}_{2\to2}
\le Cr$, together with the Hilbert bounds for $H$ and $\nabla H$,
gives, after shrinking $r$,
\begin{equation}\label{nm-tangent-hilbert}
 \norm{v}_2\le2\norm{z}_2,\qquad
 \norm{v}_2^2+\norm{\nabla v}_2^2
 \le C\langle z,Kz\rangle.
\end{equation}
No fine Hilbert-to-pointwise bound is used here.

Differentiate \eqref{nm-chart} once more. With
$H_A=(I+HD\mathcal R(A))^{-1}H$ one obtains
$A''=-H_AD^2\mathcal F(A)[v,v]$. Moreover,
$H_A^TE(A)=m(A,B)$ by \eqref{nm-flat-identities} and
\eqref{nm-multiplier}. The exact reduced Hessian is therefore
\begin{equation}\label{nm-reduced-hessian}
 \begin{split}
 D_\xi^2[\mathcal S(A(\xi,B))][z,z]
 ={}&\langle v,D^2\mathcal S(A)v\rangle\\
 &-\langle m(A,B),D^2\mathcal F(A)[v,v]\rangle.
 \end{split}
\end{equation}
By \eqref{nm-action-hessian}, \eqref{nm-tangent-energy}, and
\eqref{nm-tangent-hilbert}, the first term is at least
$(1-Cr)\langle z,Kz\rangle$. Equations
\eqref{nm-hessian-composite}, \eqref{nm-multiplier-bounds}, and
\eqref{nm-tangent-hilbert} bound the second term in absolute value
by $C(b+r^2)\langle z,Kz\rangle$. Thus throughout a sufficiently
small real chart,
\begin{equation}\label{nm-strict-hessian}
 D_\xi^2[\mathcal S(A(\xi,B))][z,z]
 \ge\tfrac12\langle z,Kz\rangle
 \ge c\norm{z}_{2,\eta}^2.
\end{equation}

\begin{theorem}[Uniform local minimum on the fixed flat gauge slice]
\label{nm-theorem}
For sufficiently small real coarse $B$, the analytic field $A(B)$
constructed above is the unique stationary point in a fixed small
chart of
\[
 \mathcal F(A)=B,\qquad RD^TA=0.
\]
It is a strict local minimum of the normalized Wilson action on
that chart. Its radius, the bounds \eqref{nm-quadratic-remainder},
the analytic derivative estimates \eqref{nm-cauchy}, and the reduced
Hessian bound \eqref{nm-strict-hessian} are independent of
refinement and compatible periodic volume.
\end{theorem}

\begin{proof}
The contraction supplies the unique small stationary field. The
converse chart following \eqref{nm-chart} includes every small
feasible field in the fixed gauge slice. Integrating the positive
Hessian \eqref{nm-strict-hessian} along a line segment in the real
$\xi$-chart proves strict local minimality and local uniqueness.
\end{proof}

\subsection{Scope and normalization}
\label{nm-scope}

The statement concerns the corrected fixed-reference constraint
and a fixed flat gauge condition. A separate gauge-orbit argument
is required to transfer it to an unrestricted constrained minimum
or to a constraint formulated with ordinary block averages. It
does not assert a minimum among large fields or a nonflat source
gauge identification.

In dimension four the coupling-free action used here agrees with
the deposited Wilson normalization: Raw I:464--479 puts $g^{-2}$
outside the sum of $1-\operatorname{Re}\tau(P)$, and
I:12756--12760 specifies the physical trace pairing. The source
left-exponential Hessian at a nonzero background is a different
coordinate Hessian from the ordinary-potential Hessian used here;
the latter has been derived directly from the same action.

On the infinite lattice the same uniformly local maps can define
bounded stationary fields, but bounded fields generally have
infinite total action. Without further finite-energy hypotheses,
one must formulate stationarity and admissible local or decaying
comparisons rather than a minimum of that total action. This
module makes its minimum claim on finite periodic regulators.

% Research proof, 16 October 2026. Two independent internal audits passed.
% No preamble. Label prefix ng-. Uses \norm and standard theorem environments.
\section{Constraint-preserving gauges and coverage of the flat slice}
\label{ng-section}

This module concerns finite compatible periodic regulators, fixed
odd blocking ratio $L$, refinement $M=L^j$, fine spacing $\eta=M^{-1}$,
and final spacing one. Fix $0<\alpha<1$. All constants below are
independent of refinement and volume. We use the corrected nonlinear
map with fixed identity reference, $\mathcal F(A)=\mathcal F^{(j)}(A)$,
and its ordinary potential differential $Q_A=D\mathcal F(A)$.
This is not renamed as the separately based source $Q(U)$.
The flat scalar projection denoted $R$ earlier is denoted $R_0$
here. The potential norms $N_\alpha$, $\mathcal N_\alpha$ and
$\mathcal X_\alpha$ refer to the same complete physical norm,
up to the fixed collar equivalences stated below.

Let $\mathcal X_\alpha$ be the complete-field norm consisting of
the bounded potential, its bounded physical first differences and
the local $\alpha$-H\"older seminorms of those differences. Let
$\mathcal Y_\alpha$ be the scalar $C^\alpha$ norm, including the
supremum, and $\mathcal Z_\alpha$ the scalar $C^{2,\alpha}$ norm,
including all physical differences through order two and the
$\alpha$-H\"older members at order two. Fixed collar enlargements
give equivalent norms. At fixed refinement these are Banach spaces.
The estimates below are uniform in that refinement. For distances
above a fixed collar size, supremum norms supply the corresponding
global increments. No field is sharply cut off before differentiating.

The purpose is to upgrade a local minimum on the fixed flat gauge
slice to a local minimum of the same nonlinear constraint, modulo
constraint-preserving gauge transformations. This requires a new
gauge argument; strict minimality on the slice alone is insufficient.

\subsection{The induced coarse gauge of the corrected average}

For the ordinary group average $M$ and one-step endpoint mean
$a_y(U)$, use the exact identities already checked in Notebook 14:
\[
 C(U)(c)=a_{c_-}(U)^{-1}M(U)(c)a_{c_+}(U),\qquad
 M(U^g)=M(U)^{g|_{\rm coarse}},
\]
where $U^g(e)=g(e_-)U(e)g(e_+)^{-1}$. Direct multiplication gives
\begin{equation}\label{ng-induced-group-gauge}
 C(U^g)=C(U)^{k(U,g)},\qquad
 k_y(U,g)=a_y(U^g)^{-1}g(y)a_y(U).
\end{equation}
All factors are on the common small-field branch. The equality is
exact, not a linearization. In particular $k(U,1)=1$.

To compute its differential, set
\[
 P_x=U(\Gamma_{yx})=e^{iY_x},\quad
 \overline Y_y=\sum_xL^{-d}Y_x,\quad
 a_y(U)=e^{i\overline Y_y},\qquad
 \mathcal J_X Z=\int_0^1\operatorname{Ad}_{e^{itX}}Z\,dt.
\]
For $g(s)=e^{is\lambda}$ the contour left increment is
$\lambda(y)-\operatorname{Ad}_{P_x}\lambda(x)$.
The logarithm and exponential differentials therefore give the
following explicit coarse gauge operator:
\begin{equation}\label{ng-one-step-gauge}
 \begin{split}
 T_U\lambda(y)=\operatorname{Ad}_{a_y(U)^{-1}}
 \bigg[\lambda(y)-\mathcal J_{\overline Y_y}
   \sum_xL^{-d}\mathcal J_{Y_x}^{-1}
       \bigl(\lambda(y)-\operatorname{Ad}_{P_x}\lambda(x)\bigr)
 \bigg].
 \end{split}
\end{equation}
Thus $(1/i)\partial_s k_y(U,e^{is\lambda})|_0=T_U\lambda(y)$.
The gauge parameter has angle units; no factor $h$ is inserted
into $g=e^{i\lambda}$. At $U=1$ this formula reduces exactly to
the normalized scalar block average $Q_{0,h}$.

For $U=e^{ihA}$, every coefficient in
\eqref{ng-one-step-gauge} is a finite local analytic function of
$hA$. Its difference from its value at zero is $O(h\norm{A}_\infty)$.
Finite one-step row incidence and physical-adjoint column incidence
then give
\begin{equation}\label{ng-one-step-smallness}
 \norm{T_U-Q_{0,h}}_{\rm av}\le C_Lh\norm{A}_\infty,
 \qquad \norm{T_U}_{\rm av}\le1+C_Lh\norm{A}_\infty.
\end{equation}
Here the absolute Schur norm uses physical measures and includes
the bounded-field norm of the transpose. Its column factor is the
fixed volume ratio $L^d$. An evaluation at the block center in
\eqref{ng-one-step-gauge} consequently costs a fixed $L$ constant,
not a power of the full refinement. The flat scalar averages and
their physical adjoints have row sums one.

Let $C_t$ be the corrected iterates of $U=e^{i\eta A}$, and let
$g_0=g$, $g_{t+1}=k(C_t,g_t)$. Induction in
\eqref{ng-induced-group-gauge} gives the exact multilevel gauge
law $C^{(j)}(U^g)=C^{(j)}(U)^{g_j}$. Differentiating at $g=1$
gives the scalar operator
\begin{equation}\label{ng-composite-gauge}
 \mathcal G_A=T_{C_{j-1}}\cdots T_{C_0},\qquad
 \mathcal G_0=Q_0.
\end{equation}
These backgrounds are the corrected iterates, unlike the ordinary
reference trajectory used for the source $Q(U)$. The proven bound
$C_t=e^{ih_tA_t}$, $\norm{A_t}_\infty\le2\norm{A}_\infty$,
and $\sum_{t<j}h_t\le1/(L-1)$ imply
\begin{equation}\label{ng-composite-smallness}
 \norm{\mathcal G_A-Q_0}_{\rm av}
       \le e^{C\norm{A}_\infty/(L-1)}-1.
\end{equation}
Its coefficients have fixed final physical range. The same estimate,
up to a fixed constant, holds in weighted physical-cell norms.
The operator is holomorphic in the common bounded-field ball;
local analytic coefficients and finite range first give holomorphy
at each refinement, and the displayed uniform bounds give uniform
Cauchy bounds on smaller balls.

Define the ordinary-coordinate fine gauge generator and its
coarse analogue by
\begin{equation}\label{ng-generators}
 X_A=\mathcal J_{\eta A}^{-1}D_U,
 \qquad U=e^{i\eta A},\quad
 D_U\lambda(e)=\eta^{-1}
       (\operatorname{Ad}_{U_e}\lambda(e_+)-\lambda(e_-)).
\end{equation}
A gauge $e^{is\lambda}$ has ordinary potential velocity
$-X_A\lambda$. At the output, with $B=\mathcal F(A)$ and
$V=e^{iB}$, differentiation of the multilevel group law yields
\begin{equation}\label{ng-gauge-intertwining}
 Q_AX_A=\mathcal J_B^{-1}D_V\mathcal G_A.
\end{equation}
The output spacing is one. In particular $X_A\lambda$ is tangent
to the nonlinear constraint whenever $\mathcal G_A\lambda=0$.

\subsection{Smooth scalar tools for the flat gauge}

Retain $\mathcal L=D^TD$, the scalar flat projection $R_0$, and
$Z$ of the previous flat notebooks. Put
\begin{equation}\label{ng-T-definition}
 T=G'R_0=Z\mathcal L.
\end{equation}
The exact identities give
\begin{equation}\label{ng-T-identities}
 Q_0T=0,\qquad R_0\mathcal L T=R_0.
\end{equation}
The following scalar estimates are needed:
\begin{equation}\label{ng-scalar-bounds}
 \norm{R_0f}_{\mathcal Y_\alpha}\le C\norm{f}_{\mathcal Y_\alpha},
 \qquad
 \norm{Tf}_{\mathcal Z_\alpha}\le C\norm{f}_{\mathcal Y_\alpha}.
\end{equation}
Here is a direct justification from the already proved canceled
aliases. With $A_0=a_0a_0^-$ and
$D_{s,2}=A_0+P_0^2\sum_{r\ne0}a_ra_r^-/P_r^2$, the high block of
$T=Z\mathcal L$ is
\[
 T_{rs}=\frac{\delta_{rs}}{P_r}
       -\frac{P_0^2a_r^-a_s}{P_r^2P_sD_{s,2}},\qquad r,s\ne0.
\]
Its mixed entries are
$T_{0s}=-a_0^-a_s/(P_sD_{s,2})$ and
$T_{r0}=-P_0a_r^-a_0/(P_r^2D_{s,2})$, while
$T_{00}=P_0R_{s,2}/D_{s,2}$.
Subtract $(\mathcal L+1)^{-1}$. The high diagonal remainder has
heat factor $1-e^{-t}$, making its two-output-difference
$\alpha$-H\"older estimate integrable. All high rank output
vectors are $a_r^-/P_r^2$, which are absolutely summable after
$2+\alpha$ differences; input rows are in alias $\ell^2$.
The scalar massive part maps $C^\alpha$ to $C^{2,\alpha}$ by
zero-row heat cancellation and the time split at the square of
the output separation. This is the same second/third heat
derivative calculation used for the right-divergence Schauder
theorem, now with both derivatives on the output.

Also $R_0=I-P_0^{\rm gauge}$, where the rank-one Bloch projection
has output vector
\[
 \chi_0=a_0^-,\qquad \chi_r=P_0a_r^-/P_r\quad(r\ne0),
 \qquad P_0^{\rm gauge}=\chi\chi^-/D_{s,2}.
\]
It maps bounded fine fields to $C^{1,\beta}$ for every fixed
$\beta<1$, by absolute output alias sums and input $\ell^2$.
This proves the first bound in \eqref{ng-scalar-bounds} for fixed
$\alpha<1$. Periodic reconstruction and contour shifting can
be applied to these complete operators; hence the same estimates
have exponential cell locality and weighted versions. These
arguments do not claim a second-difference estimate on arbitrary
bounded forcing for the massive inverse.

We also require a sufficiently smooth scalar right inverse of
$Q_0$. Define $H_b$ by the canceled scalar biharmonic symbols
\begin{equation}\label{ng-biharmonic-extension}
 (H_b)_0=\frac{a_0^-}{D_{s,2}},\qquad
 (H_b)_r=\frac{P_0^2a_r^-}{P_r^2D_{s,2}}\quad(r\ne0).
\end{equation}
Then
\begin{equation}\label{ng-biharmonic-bounds}
 Q_0H_b=I,\qquad
 \norm{H_bc}_{\mathcal Z_\alpha}\le C\norm{c}_\infty.
\end{equation}
On nonzero real fibers this is exactly
$\mathcal L^{-2}Q_0^T(Q_0\mathcal L^{-2}Q_0^T)^{-1}$.
The canceled formula supplies the constant-mode continuation, so
the complete operator has the periodic gluing needed above.
The identity follows by multiplying the scalar row $a_r$ into
\eqref{ng-biharmonic-extension}. The bound follows from the
same $a_r/P_r^2$ absolute alias sums; these even allow $3+\beta$
differences for $\beta<1$. Central entries are bounded and analytic
through $P_0=0$. This is a separately constructed scalar extension,
not the less smooth extension $G'^2Q_0^T(Q_0G'^2Q_0^T)^{-1}$.

\subsection{An invertible gauge-fixing differential}

Set $F_A=\mathcal G_A-Q_0$. On a sufficiently small field ball,
\begin{equation}\label{ng-scalar-chart}
 T_A^s=(I+H_bF_A)^{-1}
\end{equation}
is a uniformly bounded holomorphic operator on $\mathcal Z_\alpha$.
Indeed $F_A:\mathcal Z_\alpha\to\ell^\infty$ has norm
$O(\norm{A}_\infty)$ and $H_b$ satisfies
\eqref{ng-biharmonic-bounds}. The inverse differs from identity
by the same small order. If $Q_0\psi=0$ and
$\lambda=T_A^s\psi$, applying $Q_0$ to
$\lambda+H_bF_A\lambda=\psi$ gives
\begin{equation}\label{ng-scalar-kernel}
 \mathcal G_A\lambda=0.
\end{equation}

There is a uniform estimate
\begin{equation}\label{ng-generator-bound}
 \norm{(X_A-D)\lambda}_{\mathcal X_\alpha}
       \le C\norm{A}_{\mathcal X_\alpha}
                        \norm{\lambda}_{\mathcal Z_\alpha}.
\end{equation}
To see the physical scaling, write
$D_U=D+E_s(A)$ with
$E_s(A)\lambda=\eta^{-1}(\operatorname{Ad}_{e^{i\eta A}}-I)
\lambda(\cdot+\eta e_\mu)$.
The coefficient in this expression extends analytically to
$\eta=0$ and has $C^{1,\alpha}$ norm at most
$C\norm{A}_{\mathcal X_\alpha}$.
Also $\mathcal J_{\eta A}^{-1}-I$ has that norm at most
$C\eta\norm{A}_{\mathcal X_\alpha}$.
Discrete product and shift estimates then prove
\eqref{ng-generator-bound}, using two derivatives of $\lambda$.
Every derivative of a coefficient is paid by a derivative of $A$;
there is no inverse-mesh loss. These maps are holomorphic in their
indicated Banach norms on a common smaller ball.

The operator $R_0D^T:\mathcal X_\alpha\to\mathcal Y_\alpha$
is uniformly bounded. Let
$\mathcal R_\alpha=\{f\in\mathcal Y_\alpha:R_0f=f\}$,
a closed subspace because of \eqref{ng-scalar-bounds}. Define on it
\begin{equation}\label{ng-fixing-map}
 \mathcal M_A=R_0D^TX_AT_A^sT.
\end{equation}
By \eqref{ng-T-identities}, $\mathcal M_0=I$ on
$\mathcal R_\alpha$. The preceding estimates give
\begin{equation}\label{ng-fixing-smallness}
 \norm{\mathcal M_A-I}_{\mathcal R_\alpha\to\mathcal R_\alpha}
       \le C\norm{A}_{\mathcal X_\alpha}.
\end{equation}
Its range is in $\mathcal R_\alpha$ because it begins with $R_0$.
For a fixed small radius its Neumann inverse is uniformly bounded
and holomorphic. Define the scalar correction operator
\begin{equation}\label{ng-correction}
 \mathcal L_A=T_A^sT\mathcal M_A^{-1}R_0D^T.
\end{equation}
It satisfies the exact identities
\begin{equation}\label{ng-correction-identities}
 \mathcal G_A\mathcal L_Av=0,\qquad
 R_0D^TX_A\mathcal L_Av=R_0D^Tv,
\end{equation}
and the stronger norm estimate
\begin{equation}\label{ng-correction-bound}
 \norm{\mathcal L_Av}_{\mathcal Z_\alpha}
       \le C\norm{R_0D^Tv}_{\mathcal Y_\alpha}.
\end{equation}

Consequently any constrained tangent $Q_Av=0$ decomposes as
\begin{equation}\label{ng-tangent-decomposition}
 v=\bigl(v-X_A\mathcal L_Av\bigr)+X_A\mathcal L_Av,
\end{equation}
where the first term is both constrained and in the fixed flat
gauge, and the second is a constraint-preserving gauge tangent.
The first claim uses \eqref{ng-gauge-intertwining} and
\eqref{ng-correction-identities}. Its tangent space is exactly the
image of $X=\ker Q\cap\ker R_0D^T$ under the constraint-chart
derivative $T_A=(I+H D\mathcal R(A))^{-1}$: for a constrained
fixed-gauge tangent $v$, the vector
$z=(I+H D\mathcal R(A))v$ satisfies $Qz=0$ and $R_0D^Tz=0$.
A stationary Wilson field on
the fixed slice is therefore stationary on the full constraint:
the action derivative vanishes on the first term by slice
stationarity and on the second by exact gauge invariance.

\subsection{A constraint-preserving flow into the slice}

For real fields use the differential equation
\begin{equation}\label{ng-flow}
 \frac{dA(t)}{dt}=-X_{A(t)}\mathcal L_{A(t)}A(t),
 \qquad A(0)=A^{\rm in}.
\end{equation}
The vector field is locally Lipschitz on the small
$\mathcal X_\alpha$ ball, uniformly in refinement. Its negative
sign matches the actual gauge velocity in
\eqref{ng-generators}. Along any solution in that ball,
\begin{equation}\label{ng-flow-identities}
 R_0D^TA(t)=e^{-t}R_0D^TA^{\rm in},\qquad
 \mathcal F(A(t))=\mathcal F(A^{\rm in}),\qquad
 \mathcal S_\eta(e^{i\eta A(t)})
             =\mathcal S_\eta(e^{i\eta A^{\rm in}}).
\end{equation}
The first identity follows directly from
\eqref{ng-correction-identities}; the second from
\eqref{ng-gauge-intertwining}; the third from gauge invariance.
All parameter operators preserve real fields. Thus the flow
integrates unitary gauge transformations on the finite regulator.

By \eqref{ng-correction-bound} and the uniform bound for $X_A$,
\begin{equation}\label{ng-flow-speed}
 \norm{A'(t)}_{\mathcal X_\alpha}
    \le Ce^{-t}\norm{R_0D^TA^{\rm in}}_{\mathcal Y_\alpha}.
\end{equation}
Choose the initial ball sufficiently small that its radius plus
this total displacement lies within the established operator
domain. Standard continuation then keeps the solution in that
domain for every $t\ge0$. The integrable speed proves convergence
in $\mathcal X_\alpha$ to $A^\flat$, with
\begin{equation}\label{ng-flow-limit}
 R_0D^TA^\flat=0,\quad
 \mathcal F(A^\flat)=\mathcal F(A^{\rm in}),\quad
 \mathcal S_\eta(e^{i\eta A^\flat})
          =\mathcal S_\eta(e^{i\eta A^{\rm in}}),\quad
 \norm{A^\flat-A^{\rm in}}_{\mathcal X_\alpha}
       \le C\norm{R_0D^TA^{\rm in}}_{\mathcal Y_\alpha}.
\end{equation}
The gauge parameter is integrable even in $\mathcal Z_\alpha$,
so the gauge transformations themselves converge uniformly to a
unitary limit. Hence the limiting field remains in the original
gauge orbit, not merely at the same action value.

For a fixed slice field $A_*$ the same construction can be made
in a neighborhood of $A_*$. Since $R_0D^TA_*=0$,
the total displacement is bounded by
$C\norm{A^{\rm in}-A_*}_{\mathcal X_\alpha}$.
Choosing this neighborhood inside the strict slice-minimality
neighborhood ensures that $A^\flat$ belongs to it. Therefore,
if $A_*$ is a strict local minimum on
\[
 \mathcal F(A)=B,\qquad R_0D^TA=0,
\]
then it is a local minimum of $\mathcal F(A)=B$ without imposing
the flat gauge condition. Equality in a sufficiently small
neighborhood implies $A^\flat=A_*$ and hence that the initial
field is gauge equivalent to $A_*$. This is local minimality
modulo the constraint-preserving gauge orbit, not an absolute
global minimum among all lattice configurations.

The argument supplies uniform local gauge coverage for the stated
global small-field class. It does not cover arbitrary plaquette-small
gauge orbits, changing retained domains, topologically nontrivial
bundles, or a continuum quantum measure. No finite total action is
assigned to arbitrary bounded infinite-lattice fields here.

% No preamble. Notebook 15 candidate module, ot- labels.
\section{Transfer to the ordinary homogeneous averaging constraint}
\label{ot-section}

The preceding minimizer was constructed for the corrected
fixed-reference constraint $C^{(j)}(U)=D=e^{iB}$. The ordinary
constraint is $M^{(j)}(U)=D$, with $M^{(j)}$ the homogeneous
iteration of the ordinary group average. These two equalities
are distinct. This section proves a transfer between their small
field minimizers using actual gauge transformations. All fields
are on finite compatible periodic regulators; final spacing is
one. The radii and constants below are independent of refinement
and volume. The bare Wilson action is gauge invariant.

\subsection{An exact multilevel gauge identity}

Let $C_t=C^{(t)}(U)$ and $b_0(U)=1$. Notebook 14 proves
\begin{equation}\label{ot-cumulative}
 b_{t+1}(U)=b_t(U)|_{t+1}\,a(C_t),\qquad
 M^{(t)}(U)=C_t^{\,b_t(U)}.
\end{equation}
Here restriction is to the next coarse sites, and
$a(C_t)$ is the endpoint mean at that step. In particular,
on a fixed sufficiently small bounded-potential ball,
\begin{equation}\label{ot-b-bound}
 \norm{b_j(U)-1}_\infty+\norm{b_j(U)^{-1}-1}_\infty
       \le C\norm{A}_\infty,\qquad U=e^{i\eta A}.
\end{equation}
These group-valued maps are holomorphic with uniform bounds
on a common complex ball. This follows either directly from
the product in \eqref{ot-cumulative} and the summability of
the physical spacings, or from the established corrected
trajectory bounds. Only a coarse supremum bound is needed here.

For a fine gauge $g$, let $k_0=g$ and let $k_t$ be its induced
gauge under the corrected iteration, so that
$C^{(t)}(U^g)=C_t^{k_t}$. The one-step formula gives
\[
 k_{t+1}=a(C_t^{k_t})^{-1}\,k_t|_{t+1}\,a(C_t).
\]
Multiplication with \eqref{ot-cumulative}, applied also to
$U^g$, proves by induction
\begin{equation}\label{ot-exact-induced}
 b_t(U^g)k_t=g|_t\,b_t(U),\qquad
 k_j(U,g)=b_j(U^g)^{-1}\,g|_j\,b_j(U).
\end{equation}
Indeed the adjacent factors $a(C_t^{k_t})$ and their inverses
cancel before the inductive identity is used. No multiplicative
cocycle for the logarithmic endpoint average is assumed.

\subsection{Uniform analytic action of smooth fine gauges}

Write $\mathcal X_\alpha=C^{1,\alpha}$ for the complete
potential norm and $\mathcal Z_\alpha=C^{2,\alpha}$ for scalar
gauge parameters, as in the preceding section. There is a
uniform analytic map
\begin{equation}\label{ot-gauge-action}
 \mathscr A_\eta(A,\lambda)(e)
  =\frac1{i\eta}\log\bigl[
      e^{i\lambda(e_-)}e^{i\eta A(e)}e^{-i\lambda(e_+)}\bigr]
\end{equation}
from a fixed small ball in
$\mathcal X_\alpha\times\mathcal Z_\alpha$ to
$\mathcal X_\alpha$, and
\begin{equation}\label{ot-gauge-action-bound}
 \norm{\mathscr A_\eta(A,\lambda)}_{\mathcal X_\alpha}
       \le C\bigl(\norm{A}_{\mathcal X_\alpha}
                       +\norm{\lambda}_{\mathcal Z_\alpha}\bigr).
\end{equation}
To justify the apparent division by $\eta$, put
$\ell=\lambda(e_-)$ and $v=\nabla_\mu\lambda(e_-)$.
The group element in \eqref{ot-gauge-action} equals
\[
 e^{i\eta\operatorname{Ad}_{e^{i\ell}}A(e)}
        e^{i\ell}e^{-i(\ell+\eta v)}.
\]
For every small $(A,\ell,v)$ this analytic product is one
at $\eta=0$. Its logarithm is therefore divisible by $\eta$
as an analytic function, with a uniform extension for
$0\le\eta\le1$. All its finite-dimensional derivatives are
uniform on a smaller fixed complex jet ball. Discrete product,
shift and analytic-composition estimates applied to
$(A,\lambda,\nabla\lambda)$ prove
\eqref{ot-gauge-action-bound} and holomorphy in the indicated
Banach norms. A first potential difference uses at most two
scalar differences. The same argument applies to its
$\alpha$-H\"older seminorm, so no inverse-mesh loss occurs.
For real Hermitian traceless arguments the output remains
Hermitian traceless on the common logarithm branch.

\subsection{Prescribing the induced coarse gauge}

Use the smooth block-mean right inverse $H_b$ constructed in
\eqref{ng-biharmonic-extension}--\eqref{ng-biharmonic-bounds}.
It is a scalar operator tensored with the gauge Lie algebra.
For bounded coarse $\theta$ define
\begin{equation}\label{ot-K-definition}
 \mathcal K(A,\theta)=\frac1i\log
      k_j\bigl(e^{i\eta A},e^{iH_b\theta}\bigr).
\end{equation}
The exact formula \eqref{ot-exact-induced}, the bound
\eqref{ot-b-bound}, the gauge-action estimate and
$H_b:\ell^\infty\to\mathcal Z_\alpha$ show that
$\mathcal K$ is uniformly bounded and holomorphic on a
fixed ball in $\mathcal X_\alpha\times\ell^\infty$.
Evaluation of $H_b\theta$ at coarse sites is bounded by
its supremum norm. Thus this step does not presume that
$H_b$ interpolates point values.

At $\theta=0$ the induced gauge is one, so
$\mathcal K(A,0)=0$. At the joint origin its differential is
\begin{equation}\label{ot-K-linear}
 D_\theta\mathcal K(0,0)=\mathcal G_0H_b=Q_0H_b=I.
\end{equation}
Uniform Cauchy estimates on a smaller ball consequently give
\begin{equation}\label{ot-K-near-identity}
 \norm{D_\theta\mathcal K(A,\theta)-I}_{\infty\to\infty}
       \le C\bigl(\norm{A}_{\mathcal X_\alpha}
                                    +\norm{\theta}_\infty\bigr).
\end{equation}
For small $A$ and any sufficiently small bounded coarse
$\omega$, the contraction
\[
 \theta=\omega-\bigl[\mathcal K(A,\theta)-\theta\bigr]
\]
therefore has a unique solution in a fixed small ball.
The solution satisfies $\norm{\theta}_\infty\le2\norm{\omega}_\infty$
and depends holomorphically on $(A,\omega)$, with uniform
Cauchy bounds on smaller balls. It is real for real inputs.
Thus any small prescribed coarse gauge $e^{i\omega}$ is
realized as the induced corrected gauge of a fine gauge
$e^{iH_b\theta}$ of uniformly controlled smoothness.

Suppose now that $U=e^{i\eta A}$ satisfies
$M^{(j)}(U)=D$. By \eqref{ot-cumulative},
$C^{(j)}(U)^{b_j(U)}=D$. Prescribe
\[
 \omega=\frac1i\log b_j(U),\qquad k_j(U,g)=b_j(U).
\]
The preceding construction produces a fine gauge $g$ such that
\begin{equation}\label{ot-ordinary-to-corrected}
 C^{(j)}(U^g)=D,\qquad
 \mathcal S_\eta(U^g)=\mathcal S_\eta(U),\qquad
 \norm{\mathscr A_\eta(A,\log(g)/i)}_{\mathcal X_\alpha}
                     \le C\norm{A}_{\mathcal X_\alpha}.
\end{equation}
Here $\log(g)/i=H_b\theta$ on the chosen branch. This map
takes sufficiently small ordinary-feasible fields into
uniformly small corrected-feasible fields.

\subsection{An ordinary-feasible representative of the minimizer}

Let $U_*(B)=e^{i\eta A_*(B)}$ be the corrected minimizer
constructed above, so $C^{(j)}(U_*)=D=e^{iB}$ and
$\norm{A_*(B)}_{\mathcal X_\alpha}\le C\norm{B}_\infty$.
We need a smooth fine gauge with specified \emph{point values}
at the final coarse sites. For this purpose construct a
different operator $E_{\rm pt}$ as follows. Choose a fixed
smooth scalar bump $\varphi$ supported in
$(-1/2,1/2)^d$ and satisfying $\varphi(0)=1$. On the
periodic coarse lattice put
\begin{equation}\label{ot-point-extension}
 (E_{\rm pt}\omega)(x)=\sum_{n\in\mathbb Z^d}
                  \varphi(x-n)\omega(n),
\end{equation}
with the coefficients extended periodically; evaluate the
result at fine sites. This is well defined on every
compatible torus, including the smallest coarse periods.
Its derivatives through order three are bounded by
$C\norm{\omega}_\infty$, since the translates have uniformly
bounded overlap. Sampling and the integral formula for
finite differences give, uniformly in $\eta$,
\begin{equation}\label{ot-point-bounds}
 E_{\rm pt}\omega|_{\rm coarse}=\omega,\qquad
 \norm{E_{\rm pt}\omega}_{\mathcal Z_\alpha}
                          \le C\norm{\omega}_\infty.
\end{equation}
Set
\begin{equation}\label{ot-ordinary-representative}
 \omega_*(B)=\frac1i\log b_j(U_*(B))^{-1},\quad
 h_*(B)=e^{iE_{\rm pt}\omega_*(B)},\quad
 U_{\rm ord}(B)=U_*(B)^{h_*(B)}.
\end{equation}
Ordinary gauge covariance and $h_*|_{\rm coarse}=b_j(U_*)^{-1}$
give
\begin{equation}\label{ot-ordinary-exact}
 M^{(j)}(U_{\rm ord}(B))=D,\qquad
 \mathcal S_\eta(U_{\rm ord}(B))=\mathcal S_\eta(U_*(B)).
\end{equation}
Indeed with the convention $V^g=g_-Vg_+^{-1}$,
$(V^b)^h=V^{hb}$, so the two coarse gauges cancel in the
required order. The ordinary potential
$A_{\rm ord}(B)=\log U_{\rm ord}(B)/(i\eta)$ is holomorphic
on a uniform coarse supremum ball and obeys
\begin{equation}\label{ot-ordinary-bound}
 \norm{A_{\rm ord}(B)}_{\mathcal X_\alpha}
                              \le C\norm{B}_\infty.
\end{equation}

\begin{theorem}[Ordinary homogeneous constrained minimum]
\label{ot-minimum}
Fix $0<\alpha<1$, a fixed odd blocking ratio $L\ge3$,
and the compact matrix gauge group and trace normalization
used in the Wilson calculation. There are $b_0,r_0>0$,
independent of the homogeneous refinement $M=L^j$ and
finite compatible periodic volume, with the following
property. For every real coarse potential
$\norm{B}_\infty<b_0$, the field $U_{\rm ord}(B)$ in
\eqref{ot-ordinary-representative} satisfies the ordinary
constraint $M^{(j)}(U)=e^{iB}$ and minimizes the bare Wilson
action among all fields $U=e^{i\eta A}$ satisfying that
constraint and $\norm{A}_{\mathcal X_\alpha}<r_0$.
Equality in this small field class implies gauge equivalence
to $U_{\rm ord}(B)$. Its ordinary potential has the uniform
bound \eqref{ot-ordinary-bound} and a holomorphic extension
to a fixed complex coarse-data ball.
\end{theorem}

\begin{proof}
The reduced-Hessian proof gives strict convexity of the
corrected fixed-slice action on a fixed ball of the linear
coordinate $\xi\in X$, for all sufficiently small $B$.
In particular the corrected minimizer is the unique minimum
on that ball. The constraint chart covers every sufficiently
small corrected-feasible fixed-slice field, and
$\xi=(I-HQ)A$ there; this coordinate is uniformly small in
$\mathcal X_\alpha$ since $Q$ is bounded into coarse
supremum norm and $H$ maps that norm into
$\mathcal X_\alpha$.

Given an ordinary-feasible field with
$\norm{A}_{\mathcal X_\alpha}<r_0$, first apply
\eqref{ot-ordinary-to-corrected}, and then the
constraint-preserving flow into the flat slice. Both
operations preserve the action and enlarge the potential
norm by at most a fixed factor. Choose $r_0$ so small that
the resulting fixed-slice field belongs to the chart and
its $\xi$ belongs to the strict-convexity ball. The action
comparison with $U_*(B)$ now follows from that convexity.
Choose $b_0$ smaller, if necessary, so that
$\norm{A_{\rm ord}(B)}_{\mathcal X_\alpha}<r_0$ and the
minimizer coordinate belongs to the same convex ball.
Equation \eqref{ot-ordinary-exact} transfers the comparison
to $U_{\rm ord}(B)$. All radii can be chosen with fixed
strict margins and are uniform in refinement and volume.

If equality holds, strict convexity identifies the final
fixed-slice field with $U_*(B)$. Both earlier operations
are genuine gauge transformations, so the original field
is gauge equivalent to $U_{\rm ord}(B)$. Since both have
the same ordinary coarse constraint, the restriction of
a connecting gauge to the coarse sites stabilizes that
coarse field. No uniqueness of the gauge parameter is
claimed. The analyticity and norm statements follow from
\eqref{ot-b-bound}, \eqref{ot-gauge-action-bound} and the
analytic corrected minimizer.
\end{proof}

This is a minimum in a fixed small \emph{global potential}
class on each finite homogeneous regulator. It is not a
global minimum over all lattice gauge fields, nor a claim
about arbitrary local-frame or mixed-history constraints.
The source's notation $M_{\mathfrak B}(U)=D$ includes more
general geometries; only the explicitly defined ordinary
homogeneous iteration is proved here. No continuum quantum
measure or physical mass gap follows from this theorem.


\section{What this supplies to the construction program}
\label{nm15-ledger}

The new result supplies an actual nonlinear classical background
with an exact ordinary coarse constraint. The reduced stationary
equation includes both the nonlinear Wilson force and the full
curvature of the constraint manifold. It has uniform analytic
solutions, and the real action comparison survives removal of
the fixed gauge and transfer to ordinary averages. In this
homogeneous small-global-field class, the minimizer is no longer
merely an input that must be assumed.

Notebook 16 now derives a uniform local response and a controlled
expansion of the coarse classical action from these inputs.
Its weighted argument includes the derivative of the transposed
constraint map and the coarse multiplier. It proves factorial
responses with one weighted direction at every order and controls
the particular ordinary representative through its explicit gauge
formula. A fixed-exponent connected-tree bound for multiple
insertions is stronger and remains unproved.

Even a local classical coarse action would still be only part
of the required quantum renormalization step. Integration of
fluctuations brings determinants, gauge-coordinate Jacobians,
overlapping interactions and regions outside this global chart.
Those contributions must be bounded uniformly and iterated
with the coupling before a continuum probability law can be
claimed. None is replaced by the minimization theorem.

\begin{thebibliography}{9}
\bibitem{deposit15}
E.~Sussman and collaborators, \emph{Yang--Mills-LLM-experiment},
deposited raw accounts I--III at commit
\texttt{e392ef45013a581664bab76b90867170b30fe10e};
\href{https://github.com/ethanwsussman/Yang--Mills-LLM-experiment}{repository}.
Raw I:454--479 gives the Wilson normalization,
I:12756--12760 the physical bilinear pairing,
I:15471--15478 and I:15643--15648 the source quadratic
left-exponential chart, and I:43103--43107 the bare
constrained action. The ordinary constraint is displayed
at I:43098. Original journal papers are not freshly
certified by these deposited matches.
\bibitem{earlier15}
\emph{Toward a four-dimensional Yang--Mills construction},
research Notebooks 8--14 in the present workspace,
16 October 2026. Flat source-gauge extension, current and
regularity estimates; nonlinear constraint chart; nonflat
reference comparison and Wilson energy stability.
\end{thebibliography}
\end{document}
